
% CHAPTER 5

\chapter{Conclusion and Limitations}
\label{chap:conclusion}


% 5.1 MAIN FINDINGS

\section{Main Findings}

The accuracy of a variational Gaussian expansion depends not only on the number of basis functions, but also on their radial extent, exponent spacing and numerical conditioning. By comparing even-tempered Gaussian expansions with exact hydrogenic ground states, the calculations established how these factors affect energy convergence, wavefunction accuracy and the transferability of a basis between nuclear charges.

For the twelve-function hydrogen reference basis, the calculated ground-state energy was $-0.499589$ Hartree, corresponding to a relative error of $8.2170\times10^{-4}$ ($0.0822\%$). The mean radius was $1.501245$ bohr, compared with the exact value of $1.5$ bohr, and the radial probability distribution closely reproduced the exact result. These comparisons demonstrate that a relatively small Gaussian basis can describe the overall hydrogen ground state accurately. However, the agreement was not uniform across all physical quantities. The wavefunction amplitude at the nucleus was underestimated by approximately $8.5\%$, corresponding to a $16.3\%$ underestimate of the electron density. Thus, the small variational energy error did not imply equally accurate local wavefunction behaviour. This distinction is particularly important for quantities sensitive to the near-nuclear electron distribution.

The basis-size calculations showed that convergence depends strongly on the exponent distribution. Different exponent ratios produced different rates of improvement, and one sequence approached an energy-error plateau despite remaining below the chosen condition-number threshold. The plateau therefore could not be attributed to numerical instability on the available evidence alone. Increasing the number of functions does not necessarily improve the approximation efficiently if the added Gaussian widths do not address the remaining error.

The discrete parameter search demonstrated the benefit of selecting exponents rather than simply enlarging a fixed basis. Within the tested grid, five, six and eight Gaussians were the smallest bases satisfying relative energy-error targets of $10^{-2}$, $10^{-3}$ and $10^{-4}$, respectively. The selected eight-function basis achieved an error of $6.6406\times10^{-5}$, approximately 12.4 times smaller than that of the twelve-function reference basis, despite using four fewer functions. These results establish the effectiveness of exponent selection within the tested parameter ranges, rather than a globally optimal number of Gaussian functions. The selected exponent ratios also lay at the upper boundary of the search grid, indicating that the reported minimum basis sizes may depend on the chosen search limits.

The wavefunction calculations revealed a further distinction between energy accuracy and the accuracy of the state itself. For a normalised approximate ground state, the variational energy error depends quadratically on the excited-state components of the wavefunction error, with contributions weighted by their excitation energies. The energy can therefore converge more rapidly than the wavefunction-difference norm. At $N=30$, the calculated $L^2$ error was approximately $1.505\times10^{-5}$, even though the energy error had decreased substantially. The near-nuclear calculations also confirmed a limitation inherent to finite expansions of centred spherical Gaussians. Each Gaussian has zero radial derivative at the nucleus, so no finite linear combination can reproduce the nonzero electron--nucleus cusp of the exact hydrogenic ground state. The resulting local energy diverges negatively near the nucleus, despite the total variational energy remaining finite. Improving the integrated energy therefore does not remove every local error in the wavefunction.

Numerical verification supported the consistency of the implementation through independent matrix-element integration, eigenpair-residual checks, independent $L^2$ integration and nuclear-charge scaling tests. These complementary checks demonstrated the importance of assessing numerical consistency alongside agreement with an exact benchmark. The conditioning calculations also demonstrated the importance of near-linear dependence as the basis grows. Canonical orthogonalisation can remove numerically sensitive directions, but doing so reduces the available variational space. Basis size, variational accuracy and numerical reliability must therefore be considered together.

The nuclear-charge calculations provided a particularly clear example of why the physical length scale matters. Reusing the hydrogen exponents for the one-electron ion $\mathrm{Mn}^{24+}$ gave a relative energy error of approximately $54.9\%$, as the fixed basis failed to reproduce the strong contraction of the ground-state orbital. Scaling the Gaussian exponents by $Z^2$ reduced the error to the hydrogen value of approximately $0.0822\%$ and recovered the contracted radial probability distribution without increasing the basis size. This behaviour follows directly from the coordinate scaling of the non-relativistic hydrogenic Hamiltonian and was confirmed numerically.

Beyond the individual numerical results, the investigation developed a deeper understanding of how physical accuracy, basis construction and numerical reliability interact in variational calculations. This project went beyond reproducing the hydrogen ground-state energy by examining wavefunction errors, the nuclear cusp, exponent selection and near-linear dependence. Implementing parameter searches, independent numerical checks and canonical orthogonalisation also provided practical experience in identifying and addressing limitations of scientific computations, rather than relying on a single measure of convergence.

The central conclusion is that there is no universal basis size that guarantees a given level of accuracy for every physical quantity or nuclear charge. A useful Gaussian basis must cover the relevant radial length scales, meet the accuracy requirements of the intended calculation and remain sufficiently well conditioned for a reliable variational solution. Physically appropriate exponent selection can therefore be more effective than increasing the number of functions alone.


% 5.2 LIMITATIONS AND FURTHER WORK

\section{Limitations and Further Work}

The calculations were restricted to fixed-nucleus, non-relativistic, one-electron ground states represented by centred spherical Gaussian functions. This model permits direct comparison with exact hydrogenic solutions and isolates the effects of basis construction, but does not include electron--electron interactions, molecular bonding, relativistic effects or states requiring nonzero angular momentum. In particular, the exact $Z^2$ energy scaling and associated exponent scaling apply to the hydrogenic Hamiltonian used here. They cannot be assumed to hold unchanged for neutral atoms, many-electron ions or molecular systems, where additional interactions and length scales affect the wavefunctions.

A second limitation is the inability of a finite centred spherical Gaussian expansion to satisfy the exact electron--nucleus cusp condition. Although the basis can reproduce the overall energy and radial probability distribution closely, it retains an incorrect derivative at the nucleus and a divergent local energy. Future calculations could compare the Gaussian results with alternative radial functions that reproduce the cusp more naturally, or investigate explicitly cusp-corrected wavefunctions. Such comparisons would help determine how much the near-nuclear description can be improved without substantially increasing the computational cost.

The apparent plateau in the $L^2$ wavefunction error also remains unresolved. Independent integration methods agreed closely, supporting the numerical evaluation of the error, but this does not establish the cause of its convergence behaviour. A more systematic study could vary the diffuse and tight exponent limits independently, extend the basis-size range and compare different exponent distributions. Repeating selected calculations at higher numerical precision would also help distinguish basis limitations from floating-point sensitivity.

The reported minimum basis sizes are limited by the discrete parameter grid and the chosen condition-number screening criterion. Because several selected parameters occurred at the boundaries of the search range, extending the parameter domain would be particularly useful. Continuous optimisation of $\alpha$ and $\beta$ could then be compared with the discrete search. It would also be valuable to optimise against different objectives, including the $L^2$ wavefunction error, mean-radius error or near-nuclear density, rather than the energy alone. This would test whether bases selected for minimum energy error are also suitable for other physical observables, or whether different accuracy requirements favour different exponent distributions.

The numerical conditioning analysis could also be extended. More systematic tests of canonical orthogonalisation thresholds, retained basis dimensions and higher-precision eigensolvers would help identify when near-linear dependence begins to limit the reliability of the calculated energy. Such tests would also clarify how much variational flexibility can be retained while excluding numerically unreliable directions.

Finally, the same approach could be extended to excited hydrogenic states, higher angular momenta and more general atomic or molecular Hamiltonians. These calculations would test how far the conclusions about exponent coverage and conditioning remain useful beyond the spherical one-electron ground state.

The hydrogenic model therefore provides a controlled demonstration of a broader principle in scientific computing: the quality of a numerical approximation depends on how well its mathematical representation captures the physical behaviour of the system and how reliably the solution is obtained. A successful variational basis must provide the required physical accuracy while remaining numerically stable, rather than being judged by the number of functions alone.
